\clearpage
\par\bigskip\Needspace{9\baselineskip}
\pdfbookmark[2]{Zweite Methode, [4]}{ceres-second-4}
\begin{center}[4.]\end{center}
\begingroup\small
\noindent\hspace*{\parindent}[Zur Berechnung der \(Q^{(i)}\) bedient man sich der bereits früher, S. 385, angewandten Methode; man setzt
\[
4)\qquad \mathfrak{R}\mathfrak{R}=\frac{rr'\cos b}{\beta}(1+\beta\beta-2\beta\cos(D+\epsilon'-\alpha)),
\]
wo
\[
\beta=\frac{\BareRoot{(rr+r'r'+2rr'\cos b)}-\BareRoot{(rr+r'r'-2rr'\cos b)}}{\BareRoot{(rr+r'r'+2rr'\cos b)}+\BareRoot{(rr+r'r'-2rr'\cos b)}}\mathpunct{\raisebox{.45ex}{;}}
\]
sobald man also für ein Werthepaar von \(M,M'\) die Grössen \(r,r',\cos b\) und hieraus \(\beta\) berechnet hat, ergeben sich die entsprechenden Werthe der \(Q^{(i)}\) nach dem Seite 385 gegebenen Verfahren mit den nöthigen Modificationen, die sich aus der Form der Gleichungen 1 und 4 ergeben; auch \(\alpha\) rechnet sich nach bekannten Formeln. Man findet auf diese Weise für unser Beispiel mit Annahme von
\[
\begin{aligned}
\log a&=0{,}442\ 0004&\log a'&=0{,}716\ 2304\\
\log e&=8{,}894\ 7674&\log e'&=8{,}682\ 7554
\end{aligned}
\]
\[
\text{Sonnenferne}=\tilde\omega=326^{0}\,24'\,2''
\]
\SourceLeaf{411}z. B. für \(M'=\pm18^{0}\):
\[
\log r'=0{,}735\ 7693\qquad \epsilon'=-1^{0}\,36'\,46''{,}36
\]
\begingroup\footnotesize\setlength{\arraycolsep}{2pt}
\[
\begin{array}{r|c|r|c|r|r|r}
M&\log r&\log Q''&\epsilon&\alpha&\begin{gathered}\mathrm{const.}\,r\\\cos2\alpha.\,Q''\end{gathered}&\begin{gathered}\mathrm{const.}\,r\\\sin2\alpha.\,Q''\end{gathered}\\\hline
\begin{array}{@{}r@{}}+18^{0}\\-18^{0}\end{array}&\left.\vphantom{\begin{array}{c}0\\0\end{array}}\right\}0{,}473\ 4770&\begin{array}{@{}r@{}}8{,}071\ 2291\\8{,}086\ 9422\end{array}&\left.\vphantom{\begin{array}{c}0\\0\end{array}}\right\}\mp2^{0}\,32'\,25''{,}44&\begin{array}{@{}r@{}}-2^{0}\,35'\,36''{,}80\\+2\quad12\quad11{,}11\end{array}&\begin{array}{@{}r@{}}1{,}5272\\1{,}5853\end{array}&\begin{array}{@{}r@{}}-0{,}1386\\+0{,}1221\end{array}\\[3pt]
\begin{array}{@{}r@{}}+54^{0}\\-54^{0}\end{array}&\left.\vphantom{\begin{array}{c}0\\0\end{array}}\right\}0{,}463\ 1441&\begin{array}{@{}r@{}}8{,}044\ 6234\\8{,}073\ 7311\end{array}&\left.\vphantom{\begin{array}{c}0\\0\end{array}}\right\}\mp6\quad51\quad43{,}12&\begin{array}{@{}r@{}}-6\quad31\quad41{,}64\\+6\quad38\quad40{,}63\end{array}&\begin{array}{@{}r@{}}1{,}3720\\1{,}4657\end{array}&\begin{array}{@{}r@{}}-0{,}3181\\+0{,}3462\end{array}\\[3pt]
\begin{array}{@{}r@{}}+90^{0}\\-90^{0}\end{array}&\left.\vphantom{\begin{array}{c}0\\0\end{array}}\right\}0{,}444\ 6563&\begin{array}{@{}r@{}}8{,}007\ 0532\\8{,}015\ 2581\end{array}&\left.\vphantom{\begin{array}{c}0\\0\end{array}}\right\}\mp8\quad57\quad24{,}14&\begin{array}{@{}r@{}}-8\quad34\quad46{,}38\\+9\quad9\quad49{,}82\end{array}&\begin{array}{@{}r@{}}1{,}1828\\1{,}1975\end{array}&\begin{array}{@{}r@{}}-0{,}3652\\+0{,}3967\end{array}\\[3pt]
\begin{array}{@{}r@{}}+126^{0}\\-126^{0}\end{array}&\left.\vphantom{\begin{array}{c}0\\0\end{array}}\right\}0{,}423\ 4498&\begin{array}{@{}r@{}}7{,}949\ 3412\\7{,}929\ 3295\end{array}&\left.\vphantom{\begin{array}{c}0\\0\end{array}}\right\}\mp7\quad41\quad47{,}10&\begin{array}{@{}r@{}}-7\quad44\quad29{,}95\\+8\quad6\quad45{,}95\end{array}&\begin{array}{@{}r@{}}0{,}9947\\0{,}9464\end{array}&\begin{array}{@{}r@{}}-0{,}2756\\+0{,}2754\end{array}\\[3pt]
\begin{array}{@{}r@{}}+162^{0}\\-162^{0}\end{array}&\left.\vphantom{\begin{array}{c}0\\0\end{array}}\right\}0{,}408\ 6211&\begin{array}{@{}r@{}}7{,}887\ 3044\\7{,}872\ 3513\end{array}&\left.\vphantom{\begin{array}{c}0\\0\end{array}}\right\}\mp3\quad3\quad44{,}12&\begin{array}{@{}r@{}}-3\quad24\quad40{,}10\\+3\quad4\quad55{,}48\end{array}&\begin{array}{@{}r@{}}0{,}8586\\0{,}8307\end{array}&\begin{array}{@{}r@{}}-0{,}1027\\+0{,}0897\end{array}
\end{array}
\]
\endgroup
woraus man nach den obigen Interpolationsformeln erhält
\[
\begin{aligned}
\mathrm{const.}\,r\cos2\alpha.\,Q''={}&1{,}1961+0{,}3761\cos M\ .\ .\ .\ .\\
&\quad .\ .\ .\ .\ .\ .\ .\ .\ .\ .\ .\ .\ .\ .\ .\ .\\
\mathrm{const.}\,r\sin2\alpha.\,Q''={}&.\ .\ .\ .\ .\ .\ .\ .\ .\ .\ .\ .\ .\ .\ .\ .\ .\ .\ .\ .\ .\\
&\quad-0{,}3770\sin M-\ldots,
\end{aligned}
\]
wo nur die grössten Glieder berechnet sind. Dieselbe Rechnung ist auszuführen für die übrigen Werthe von \(M'\); es fand sich so mit Berücksichtigung nur der grössten Glieder
\[
\begin{array}{r|l|l}
M'&\mathrm{const.}\,r\cos2\alpha.\,Q''&\mathrm{const.}\,r\sin2\alpha.\,Q''\\\hline
\pm18^{0}&1{,}1961+0{,}3761\cos M&-0{,}3770\sin M\\
\pm54^{0}&1{,}3186+0{,}4188\cos M&-0{,}4157\sin M\\
\pm90^{0}&1{,}5632+0{,}5057\cos M&-0{,}4929\sin M\\
\pm126^{0}&1{,}8871+0{,}6230\cos M&-0{,}5953\sin M\\
\pm162^{0}&2{,}1429+0{,}7179\cos M&-0{,}6758\sin M
\end{array}
\]
und hieraus
\[
\begin{array}{r|l}
M'&\begin{aligned}&\mathrm{const.}\,r\cos2\alpha.\,Q''\sin(3D+\text{\AstroSymbols ⚳}-\AscendingNode)\\&-\mathrm{const.}\,r\sin2\alpha.\,Q''\cos(3D+\text{\AstroSymbols ⚳}-\AscendingNode)\end{aligned}\\\hline
\pm18^{0}&\begin{aligned}&1{,}1961\sin(3D+\text{\AstroSymbols ⚳}-\AscendingNode)\\&+0{,}3766\sin(3D+\text{\AstroSymbols ⚳}-\AscendingNode+M)\end{aligned}\\[3pt]
\pm54^{0}&\begin{aligned}&1{,}3186\quad\text{»}\\&+0{,}4172\quad\text{»}\end{aligned}\\[3pt]
\pm90^{0}&\begin{aligned}&1{,}5632\quad\text{»}\\&+0{,}4993\quad\text{»}\end{aligned}\\[3pt]
\pm126^{0}&\begin{aligned}&1{,}8871\quad\text{»}\\&+0{,}6092\quad\text{»}\end{aligned}\\[3pt]
\pm162^{0}&\begin{aligned}&2{,}1429\quad\text{»}\\&+0{,}6968\quad\text{»}\end{aligned}
\end{array}
\]

Für die Grösse
\[
\begin{aligned}
&\mathrm{const.}\,r\cos2\alpha.\,Q''\cos(3D+\text{\AstroSymbols ⚳}-\AscendingNode)\\
&+\mathrm{const.}\,r\sin2\alpha.\,Q''\sin(3D+\text{\AstroSymbols ⚳}-\AscendingNode)
\end{aligned}
\]
gelten dieselben Werthe, wie vorstehend, nur sind statt \(\sin(3D+\text{\AstroSymbols ⚳}-\AscendingNode)\) und \(\sin(3D+\text{\AstroSymbols ⚳}-\AscendingNode+M)\) die entsprechenden Cosinus zu setzen.

\SourceLeaf{412}Die vorigen Ausdrücke sind nun endlich mit \(r'\cos3\epsilon'\) resp. \(r'\sin3\epsilon'\) zu multipliciren; es fand sich:
\begingroup\footnotesize\setlength{\arraycolsep}{4pt}
\[
\begin{array}{r|r|r|r|r}
M'&\log r'&\epsilon'&\log r'\cos3\epsilon'&\log(\mp r'\sin3\epsilon')\\\hline
\pm18^{0}&0{,}735\ 7693&\mp1^{0}\,36'\,46''{,}36&0{,}73422&9{,}66185\\
\pm54^{0}&0{,}728\ 9695&\mp4\quad18\quad30{,}90&0{,}71782&0{,}07861\\
\pm90^{0}&0{,}717\ 2353&\mp5\quad30\quad39{,}96&0{,}69890&0{,}17143\\
\pm126^{0}&0{,}704\ 4644&\mp4\quad37\quad26{,}37&0{,}69161&0{,}08423\\
\pm162^{0}&0{,}695\ 9734&\mp1\quad48\quad31{,}26&0{,}69402&9{,}67169
\end{array}
\]
\endgroup
und hieraus
\begin{landscape}
\begingroup\footnotesize\setlength{\arraycolsep}{5pt}
\[
\begin{array}{r|l|l}
M'&r'\cos3\epsilon'\left\{\begin{aligned}&\mathrm{const.}\,r\cos2\alpha.\,Q''\sin(3D+\text{\AstroSymbols ⚳}-\AscendingNode)\\&-\mathrm{const.}\,r\sin2\alpha.\,Q''\cos(3D+\text{\AstroSymbols ⚳}-\AscendingNode)\end{aligned}\right\}&\mp r'\sin3\epsilon'\left\{\begin{aligned}&\mathrm{const.}\,r\cos2\alpha.\,Q''\cos(3D+\text{\AstroSymbols ⚳}-\AscendingNode)\\&+\mathrm{const.}\,r\sin2\alpha.\,Q''\sin(3D+\text{\AstroSymbols ⚳}-\AscendingNode)\end{aligned}\right\}\\\hline
\pm18^{0}&\begin{aligned}&6{,}4862\sin(3D+\text{\AstroSymbols ⚳}-\AscendingNode)\\&+2{,}0419\sin(3D+\text{\AstroSymbols ⚳}-\AscendingNode+M)\end{aligned}&\begin{aligned}&0{,}5491\cos(3D+\text{\AstroSymbols ⚳}-\AscendingNode)\\&+0{,}1729\cos(3D+\text{\AstroSymbols ⚳}-\AscendingNode+M\end{aligned}\\[4pt]
\pm54^{0}&\begin{aligned}&6{,}8854\quad\text{»}\\&+2{,}1787\quad\text{»}\end{aligned}&\begin{aligned}&1{,}5802\quad\text{»}\\&+0{,}4999\quad\text{»}\end{aligned}\\[4pt]
\pm90^{0}&\begin{aligned}&7{,}8147\quad\text{»}\\&+2{,}4961\quad\text{»}\end{aligned}&\begin{aligned}&2{,}3197\quad\text{»}\\&+0{,}7410\quad\text{»}\end{aligned}\\[4pt]
\pm126^{0}&\begin{aligned}&9{,}2768\quad\text{»}\\&+2{,}9945\quad\text{»}\end{aligned}&\begin{aligned}&2{,}2910\quad\text{»}\\&+0{,}7396\quad\text{»}\end{aligned}\\[4pt]
\pm162^{0}&\begin{aligned}&10{,}5930\quad\text{»}\\&+3{,}4448\quad\text{»}\end{aligned}&\begin{aligned}&1{,}0062\quad\text{»}\\&+0{,}3272\quad\text{»}\end{aligned}
\end{array}
\]
\endgroup
\end{landscape}

Wird hienach noch die Interpolation nach dem Argument \(M'\) vorgenommen, so finden sich die Hauptglieder der Entwickelung:
\[
\begin{aligned}
&r'\cos3\epsilon'\left\{\begin{aligned}&\mathrm{const.}\,r\cos2\alpha.\,Q''\sin(3D+\text{\AstroSymbols ⚳}-\AscendingNode)\\&-\mathrm{const.}\,r\sin2\alpha.\,Q''\cos(3D+\text{\AstroSymbols ⚳}-\AscendingNode)\end{aligned}\right\}=\\
&\quad(8''{,}2112-2''{,}1246\cos M')\sin(3D+\text{\AstroSymbols ⚳}-\AscendingNode)\\
&\quad+(2''{,}6312-0''{,}7243\cos M')\sin(3D+\text{\AstroSymbols ⚳}-\AscendingNode+M)
\end{aligned}
\]
\[
\begin{aligned}
&r'\sin3\epsilon'\left\{\begin{aligned}&\mathrm{const.}\,r\cos2\alpha.\,Q''\cos(3D+\text{\AstroSymbols ⚳}-\AscendingNode)\\&+\mathrm{const.}\,r\sin2\alpha.\,Q''\sin(3D+\text{\AstroSymbols ⚳}-\AscendingNode)\end{aligned}\right\}=\\
&\quad-2''{,}3726\sin M'.\cos(3D+\text{\AstroSymbols ⚳}-\AscendingNode)\\
&\quad-0''{,}7592\sin M'.\cos(3D+\text{\AstroSymbols ⚳}-\AscendingNode+M)
\end{aligned}
\]

Die Summe dieser beiden Ausdrücke gibt die Entwickelung des Gliedes 3).]
\par\endgroup
\begin{center}\rule{.30\linewidth}{.4pt}\end{center}
